One may suppose $x\ne0$; then $\Delta(x)\ne0$ since $(\varepsilon\otimes\mathrm{id})\Delta(x)=x$, where $\varepsilon$ denotes the augmentation (counit) of $A$. Write
% typo: corrupted printed coünité -> counit.
\begin{equation*}\tag{1}
\Delta(x)=\sum_{j=1}^n e_j\otimes a_j\qquad\text{with $n$ minimal},
\end{equation*}
in which case the $e_j$ (resp. $a_j$) are linearly independent. Complete $(e_1,\ldots,e_n)$ to a basis $(e_i)_{i\in I}$ of $A$ and put, for $j\in J=\{1,\ldots,n\}$,
\[
\Delta(e_j)=\sum_{i\in I}e_i\otimes b_{ij}.
\]
Applying $\Delta\otimes\mathrm{id}$ and $\mathrm{id}\otimes\Delta$ to (1), one obtains from axiom (CO 1) of 11.8.0 (see also (HA 1) in I 4.2) the equalities:
\[
\sum_{j\in J}e_j\otimes\Delta(a_j)=\sum_{\ell\in J}\Delta(e_\ell)\otimes a_\ell=\sum_{i\in I}e_i\otimes\Bigl(\sum_{\ell\in J}b_{i\ell}\otimes a_\ell\Bigr).
\]
Since the $e_i$ are linearly independent, it follows that
\begin{equation*}\tag{2}
\forall j\in J,\quad\Delta(a_j)=\sum_{\ell\in J}b_{j\ell}\otimes a_\ell.
\end{equation*}
Let $B$ then be the sub-$k$-algebra of finite type of $A$ generated by the $b_{ij}$ and the $u(b_{ij})$, for $i,j\in J$. It is clear that $u(B)=B$.

Applying $\Delta\otimes\mathrm{id}$ and $\mathrm{id}\otimes\Delta$ to (2), one again obtains from (CO 1) the equalities:
\[
\sum_{\ell\in J}\Delta(b_{j\ell})\otimes a_\ell=\sum_{i\in J}b_{ji}\otimes\Delta(a_i)=\sum_{i,\ell\in J}b_{ji}\otimes b_{i\ell}\otimes a_\ell,
\]
and since the $a_\ell$ are linearly independent, one deduces that
\begin{equation*}\tag{3}
\forall j,\ell\in J,\quad\Delta(b_{j\ell})=\sum_{i\in J}b_{ji}\otimes b_{i\ell}.
\end{equation*}
Since $\Delta\circ u=(u\times u)\circ v\circ\Delta$, where $v(a\otimes b)=b\otimes a$, one therefore also has
\begin{equation*}\tag{4}
\forall j,\ell\in J,\quad\Delta\bigl(u(b_{j\ell})\bigr)=\sum_{i\in J}u(b_{i\ell})\otimes u(b_{ji}).
\end{equation*}
Since $\Delta$ is an algebra homomorphism, one deduces from (3) and (4) that $\Delta(B)\subset B\otimes_k B$. Finally, axiom (CO 2) of 11.8.0 (see also (HA 2) in I, 4.2) shows that $a_j=\sum_{i\in I}\varepsilon(a_j)b_{ij}$ and $x=\sum_{j\in J}\varepsilon(e_j)a_j$, so that $x\in B$.
% typo?: kept as printed: epsilon(a_j), rather than epsilon(a_i), in the first sum.

\begin{proposition}{11.13}\label{VIB.11.13}
Let $k$ be a field and $G$ an affine $k$-group with algebra $A$. Then $G$ is the projective limit of an increasing filtered system of affine $k$-groups of finite type, whose transition morphisms are faithfully flat.
\end{proposition}
If $B$ and $B'$ are two subalgebras of $A$ of finite type stable under $\Delta$ and $u$, then so is the subalgebra generated by $B$ and $B'$. Thus, by Lemma 11.12, $A$ is the inductive limit of an increasing filtered family $(B_i)_{i\in I}$ of subalgebras of finite type stable under $\Delta$ and $u$. Then each $B_i$, endowed with the restriction of $u$ and the morphism $B_i\to B_i\otimes_k B_i$ obtained from $\Delta$, is a Hopf algebra, so by I 4.2 it is the algebra of an affine $k$-group $G_i$, of finite type over $k$. Finally, since $A=\varinjlim B_i$, one has $G=\varprojlim G_i$ (cf. EGA IV$_3$ 8.2.3). The transition morphisms are faithfully flat by the following lemma:

\begin{lemma}{11.14}\label{VIB.11.14}
Let $k$ be a field, $u:G\to H$ a morphism between affine $k$-groups of finite type, and $u^\natural:B\to A$\nde{The morphism $B\to A$ corresponding to $u:\Spec(A)\to\Spec(B)$ has been denoted by $u^\natural$ (instead of $u^\circ$).} the corresponding morphism of $k$-algebras. For $u$ to be faithfully flat, it is necessary and sufficient that $u^\natural$ be injective.
\end{lemma}
The condition is obviously necessary (cf. EGA 0$_{\mathrm I}$ 6.6.1). Let us show that it is sufficient. Put $N=\Ker u$. Then, by VIA, 3.3.2 and 5.4.1, $G/N$ is a $k$-group of finite type and $u$ factors as $G\xrightarrow{p}G/N\xrightarrow{v}H$, where $p$ is faithfully flat and $v$ is a closed immersion. Thus, since $H$ is an affine scheme, $G/N$ is an affine scheme and the morphism $v^\natural:B\to\OO(G/N)$ is surjective (cf. EGA I 4.2.3). Now, since $u^\natural$ is supposed injective and $u^\natural=p^\natural\circ v^\natural$, then $v^\natural$ is also injective: it is therefore an isomorphism, as is $v$, and since $p$ is faithfully flat, so is $u$.

\begin{definition}{11.15}\label{VIB.11.15}
\nde{The hypothesis that $G$ be quasi-compact has been added, and the equivalence of conditions (i) and (ii) expanded.}
Let $k$ be a field, $G$ a quasi-compact $k$-group and $V$ a $k$-vector space endowed with a $k$-linear action of $G$, hence with an $\mathscr A(G)$-comodule structure $\delta:V\to V\otimes\mathscr A(G)$, by 11.6.1. Let $v\in V$ be nonzero. The following conditions are equivalent:

(i) There is $\lambda\in\mathscr A(G)$ (necessarily unique) such that $\delta(v)=v\otimes\lambda$.

(ii) For every $k$-algebra $R$ and every $g\in G(R)$, one has $g\cdot v\in Rv$ (i.e., there is $f\in R$, necessarily unique, such that one has in $V\otimes R$ the equality $g\cdot v=v\otimes f$).

Indeed, it is clear that (i) $\Rightarrow$ (ii). Conversely, if (ii) holds and one applies it to $R=\mathscr A(G)$ and $g=\mathrm{id}_{\mathscr A(G)}$, one obtains that there is a unique $\lambda\in\mathscr A(G)$ such that $\delta(v)=v\otimes\lambda$.

If $v$ satisfies these conditions, one says that $v$ is a \emph{semi-invariant vector under $G$}, and $\lambda$ is the \emph{weight of $v$}; one will also say that ``$v$ is a semi-invariant of weight $\lambda$''.
\end{definition}
Denote by $\Delta$ the comultiplication of $\mathscr A(G)$; then the equality
\[
v\otimes\lambda\otimes\lambda=(\delta\otimes\mathrm{id})(\delta(v))=(\mathrm{id}\otimes\Delta)(\delta(v))=v\otimes\Delta(\lambda)
\]
implies that $\Delta(\lambda)=\lambda\otimes\lambda$. Consequently, $\lambda$ defines a morphism of Hopf algebras
\[
\mathscr A(\Gm{}_{,k})=k[T,T^{-1}]\to\mathscr A(G),\qquad T\mto\lambda,
\]
and hence a morphism of $k$-groups $\lambda:G\to\Gm{}_{,k}$, i.e.\ $\lambda$ is a character of $G$, called the \emph{character associated to the semi-invariant vector $v$}.

\begin{lemma}{11.16.0}\label{VIB.11.16.0}
\nde{This lemma, taken from [DG70], \S\,II.2, 3.5, has been added.}
Let $k$ be a field, $H$ an affine $k$-group, $V$ an $H$-module of dimension $n$ and $U$ a vector subspace of $V$ of dimension $d$. Consider the line $D=\bigwedge^d U\subset\bigwedge^d V$. For $U$ to be stable under $H$, it is necessary and sufficient that $D$ be so.
\end{lemma}
Since necessity is clear, let us prove sufficiency. One may suppose $d<n$. Let $(e_1,\ldots,e_d)$ be a basis of $U$; complete it to a basis $(e_1,\ldots,e_n)$ of $V$. For every $k$-algebra $R$, $V_R=V\otimes R$ is a free $R$-module and one has
\[
U_R=\{v\in V_R\mid v\wedge(e_1\wedge\cdots\wedge e_d)=0\}
\]
(for, for $i>d$, the $e_i\wedge e_1\wedge\cdots\wedge e_d$ are linearly independent in $\bigwedge_R^{d+1}V_R$). Since $H(R)$ acts on $\bigwedge_R^\bullet(V_R)$ by
\[
h(x_1\wedge\cdots\wedge x_s)=h(x_1)\wedge\cdots\wedge h(x_s),
\]
it follows that, if $H(R)$ stabilizes $D_R=Re_1\wedge\cdots\wedge e_d$, it also stabilizes $U_R$.

\begin{theorem}{11.16}\label{VIB.11.16}\textup{(Chevalley).}
Let $k$ be a field, $G$ an affine algebraic $k$-group, $H$ a closed subgroup scheme of $G$.\nde{On the one hand, recall that every subgroup scheme of $G$ is closed (1.4.2). On the other hand, the result has been stated in the usual form: ``$H$ is the stabilizer of a line in a representation of $G$'', while retaining the original formulation in terms of a sequence $a_1,\ldots,a_n$ of semi-invariants in $\mathscr A(G)$.} Then there is a finite-dimensional $G$-module $V$ and a line $D$ of $V$ such that $H=\uNorm_G(D)$, i.e.\ such that for every $k$-algebra $R$,
\[
H(R)=\{g\in G(R)\mid g(D_R)=D_R\}.
\]
In other words, there are finitely many elements $a_1,\ldots,a_n\in\mathscr A(G)$ which are semi-invariants, all of the same weight $\lambda$, for the ``right'' action of $H$ (i.e.\ $a_i(gh)=\lambda(h)a_i(g)$, for every $k$-algebra $R$ and $g\in G(R)$, $h\in H(R)$), such that $H$ is the largest closed subgroup scheme of $G$ under which the $a_i$ are semi-invariants.
\end{theorem}
Denote by $\Delta$ (resp. $\varepsilon$) the comultiplication (resp. augmentation) of $A=\mathscr A(G)$. Then $H=\Spec(A/I)$, for some ideal $I$ of $A$ contained in $\Ker\varepsilon$ and such that $\Delta(I)\subset I\otimes A+A\otimes I$. Let $B=A/I$ and $\pi$ be the projection $A\to B$. Consider the left action of $H$ on $A$ given by $(h\varphi)(g)=\varphi(gh)$; the corresponding $B$-comodule structure is given by:
\[
\overline\Delta:\quad A\xrightarrow{\Delta}A\otimes A\xrightarrow{\mathrm{id}_A\otimes\pi}A\otimes B.
\]
Then $I$ is a sub-$H$-module of $A$, since $\overline\Delta(I)\subset I\otimes B$.

On the other hand, $A$ is a $k$-algebra of finite type, hence noetherian, so $I$ has a finite set of generators $(x_1,\ldots,x_r)$. By 11.8, the $x_i$ are contained in a sub-$G$-module $V$ finite-dimensional over $k$. Then $W=V\cap I$ is a finite-dimensional $H$-module, whose dimension we shall denote by $d$. Since $V$ contains all the $x_i$, $W$ generates the ideal $I$.

Put $E=\bigwedge^d V$, let $(w_1,\ldots,w_d)$ be a basis of $W$, and let $\mathscr B=(e_0,\ldots,e_n)$ be a basis of $E$ containing the vector $e_0=w_1\wedge\cdots\wedge w_d$. The operation of $G$ on $V$ canonically determines an operation of $G$ on $E$, hence an $A$-comodule structure $\rho:E\to E\otimes A$. For $j=0,\ldots,n$, put $\rho(e_j)=\sum_{i=1}^n e_i\otimes b_{ij}$; we saw in the proof of 11.12 that then
\begin{equation*}\tag{*}
\Delta(b_{ij})=\sum_{\ell=0}^n b_{i\ell}\otimes b_{\ell j}.
\end{equation*}
% typo?: kept as printed: rho(e_j) summed from i=1, while the basis starts at e_0.

Put $a_i=b_{i0}$, i.e.\ if $(e_0^*,\ldots,e_n^*)$ is the dual basis of $\mathscr B$, the $a_i$ are the ``matrix coefficients'' $g\mto e_i^*(ge_0)$. On the other hand, the operation of $H$ on $E$ corresponds to:
\[
\overline\rho:\quad E\xrightarrow{\rho}E\otimes A\xrightarrow{\mathrm{id}_E\otimes\pi}E\otimes B.
\]
Since $W$ is stable under $H$, then $e_0$ is semi-invariant under $H$, so $\overline\rho(e_0)=e_0\otimes\pi(a_0)$ and $a_i=b_{i0}$ belongs to $I$ for $i=1,\ldots,n$. Substituting this into (*), one obtains $\overline\Delta(a_i)=a_i\otimes\pi(a_0)$ for $i=0,\ldots,n$, i.e.\ the $a_i$ are semi-invariants under $H$ of weight $\pi(a_0)$. (Moreover, replacing $E$ if necessary by the sub-$G$-module generated by $e_0$ (cf. 11.8), one may suppose that the $a_i$ are linearly independent.)

Conversely, let $H'=\Spec B'$, where $B'=A/I'$, be a closed subgroup scheme of $G$ under which each $a_i$ is semi-invariant, of weight $\lambda_i\in B'$ (this is the case, in particular, if $e_0$ is invariant under $H'$ of weight $\lambda'$). Let us show that $H'=H$. Denote by $\pi'$ the projection $A\to B'$; the hypothesis implies that
\[
a_i\otimes\lambda_i=(\mathrm{id}_A\otimes\pi')\Delta(b_{i0})=\sum_{\ell=0}^n b_{i\ell}\otimes\pi'(a_\ell),
\]
whence $\lambda_i=\pi'(a_0)$ and $a_\ell\in I'$ for $\ell=1,\ldots,n$, and therefore $e_0$ is semi-invariant under $H'$. By Lemma 11.16.0, this implies that $W$ is stable under $H'$. Since the ideal $I$ is generated by $W$, it is therefore also stable under $H'$, and hence
\[
\Delta(I)\subset I\otimes A+A\otimes I'.
\]
Since $I\subset\Ker\varepsilon$ and $(\varepsilon\otimes\mathrm{id}_A)\circ\Delta=\mathrm{id}_A$, it follows that $I\subset I'$, whence $H'\subset H$.
% typo?: kept as printed: the proof announces H'=H, but concludes H' subset H.

\begin{lemma}{11.17.0}\label{VIB.11.17.0}
\nde{This lemma, taken from the proof of thm. 5.6 of [DG70], \S\,III.3, has been added (by abuse of notation, the same letter $E$ denotes a $k$-vector space and the $k$-scheme in modules $\mathbf W(E)=\Spec S(E^*)$).}
Let $k$ be an algebraically closed field, $G$ a reduced affine algebraic $k$-group, $V$ a finite-dimensional $G$-module over $k$, and $v\in V$. Let $E$ be the vector subspace of $V$ generated by the vectors $gv$, for $g\in G(k)$. Then $E$ is the smallest sub-$G$-module of $V$ containing $v$, and thus the morphism $G\times E\to V$, $(g,x)\mto gx$ factors through $E$.
\end{lemma}
\begin{proof}
By 11.8, one knows that there is a smallest sub-$G$-module $U$ of $V$ containing $v$: if $\mu:V\to V\otimes\mathscr A(G)$ denotes the comodule structure and one writes $\mu(v)=\sum_{i=1}^n v_i\otimes a_i$ with the $a_i$ linearly independent, one has $U=\mathrm{Vect}(v_1,\ldots,v_n)$. It is clear that $U$ contains $E$, and the morphism $G\times U\to V$ factors through $U$.

Conversely, the inverse image of $E$ under the morphism $\mu_v:g\mto gv$ is a closed subset of $G$ which contains the rational points; now these are dense in $G$, since $G$ is of finite type over $k$ (cf. EGA IV$_3$, 10.4.8), so $\mu_v^{-1}(E)=G$ and thus, since $G$ is reduced, $\mu_v$ factors through $E$, whence $E=U$.
\end{proof}

\begin{theorem}{11.17}\label{VIB.11.17}\textup{(Chevalley).}
Let $k$ be a field, $G$ an affine $k$-group (not necessarily of finite type), and $N$ a closed subgroup scheme of $G$ invariant in $G$; then the (fpqc) quotient sheaf $G/N$ is representable by an affine $k$-group.\nde{For another proof of this theorem, which does not use the results of VIA, see [Ta72], Th. 5.2 (see also Remark 11.18.5).}
\end{theorem}
Suppose first $G$ of finite type. By VIA 3.2 and 5.2, the (fpqc) quotient sheaf $G/N$ is representable by a $k$-group $Q$; it is therefore a matter of showing that $Q$ is affine. The proof proceeds in several steps; suppose first $k$ algebraically closed.\nde{In what follows the original has been expanded (and the erroneous assertion $(G/N)_{\mathrm{red}}=G_{\mathrm{red}}/N_{\mathrm{red}}$ corrected), following [DG70], \S\,III.3, 5.6.}

(a) Suppose moreover $G$ reduced and connected and $N$ reduced. By 11.16, there is a $G$-module $V$ finite-dimensional over $k$, and a line $D=ke_0$ such that $\uNorm_G(D)=N$; in particular $N$ acts on $D$ through a character $\chi:N\to\Gm{}_{,k}$.

Fix $h\in N(k)$. For every $g\in G(k)$, one has $hge_0=g(g^{-1}hg)e_0=\chi(g^{-1}hg)ge_0$, so $\chi(g^{-1}hg)$ is an eigenvalue of $h$. Thus the continuous map $\phi:G(k)\to k$, $g\mto\chi(g^{-1}hg)$, takes only finitely many values, and since $G(k)$ is irreducible (for it is dense in $G$), one therefore has $\phi(g)=\phi(e)=\chi(h)$ for every $g\in G(k)$ and hence
\[
\chi(g^{-1}hg)=\chi(h),\qquad\forall g\in G(k),\ h\in N(k).
\]
Let $E$ be the vector subspace of $V$ generated by the vectors $ge_0$, for $g\in G(k)$; by Lemma 11.17.0, it is the sub-$G$-module of $V$ generated by $e_0$.

By the above, the two morphisms $N\times E\to E$, $(h,x)\mto hx$ and $(h,x)\mto\chi(h)x$, coincide on the set of rational points $(N\times E)(k)=N(k)\times E(k)$, which is dense in $N\times E$. Since $N\times E$ is reduced (and $E$ separated), these two morphisms are therefore equal, so $N$ acts on $E$ by homotheties. Consequently, $N$ is contained in the kernel $K$ of the morphism $\rho:G\to\GL(\End_k(E))$, defined by $\rho(g)(u)=gug^{-1}$, for every $g\in G(R)$ and $u\in\End_R(E\otimes R)$ ($R$ a $k$-algebra). On the other hand, if $g\in K(R)$ then $g(Re_0)=Re_0$, whence $g\in N(R)$. This shows that $N=K$. Then, by VIA 5.4.1, the morphism $G/N\to\GL(\End_k(E))$ is a closed immersion, and thus $G/N$ is affine.

(b) Suppose now $G$ and $N$ reduced, with $G$ not necessarily connected. Put $N'=N\cap G^0$; then $G^0/N'$ is affine by (a). On the other hand, $NG^0$ is an invariant subgroup of $G$ and $G/NG^0$, being a quotient of the finite constant group $G/G^0$ (cf. VIA, 5.5.1), is likewise a finite constant group. Thus $G/N$ is the direct sum of the fibers of the morphism $G/N\to G/NG^0$, all isomorphic to $NG^0/N$, hence to $G^0/N'$, by VIA, 5.3.3. Thus $G/N$ is affine.

(c) Suppose $G$ reduced, and $N$ arbitrary. The morphism $G\times N\to N$, $(g,h)\mto ghg^{-1}$, induces a morphism $(G\times N)_{\mathrm{red}}\to N_{\mathrm{red}}$; now, since $G$ is reduced and $k$ algebraically closed, one has $(G\times N)_{\mathrm{red}}=G\times N_{\mathrm{red}}$, so $N_{\mathrm{red}}$ is an invariant subgroup of $G$. (N.B. This fails when $G$ is not reduced, cf. VIA, 0.2.)

Thus, by (a), $G'=G/N_{\mathrm{red}}$ is affine. On the other hand, by VIA 5.6.1, $N'=N/N_{\mathrm{red}}$ is a finite $k$-group, so, by Theorem 4.1 of Exp. V, the quotient $G/N=G'/N'$ is affine.
% typo?: kept as printed: reference to (a), rather than (b), here.

(d) For arbitrary $G$ and $N$, the equivalence relation obtained from $G\times N\rightrightarrows G$ by the base change $G_{\mathrm{red}}\to G$ is:
\[
G_{\mathrm{red}}\times N'\rightrightarrows G_{\mathrm{red}},\qquad\text{where}\qquad N'=N\cap G_{\mathrm{red}}.
\]
Since the underlying spaces are the same (and since the quotient is the quotient ringed space), the morphism $G_{\mathrm{red}}/N'\to G/N$ is a homeomorphism. Since $G_{\mathrm{red}}/N'$ is reduced (for $p:G_{\mathrm{red}}\to G_{\mathrm{red}}/N'$ is faithfully flat), it follows that $(G/N)_{\mathrm{red}}$ is identified with $G_{\mathrm{red}}/N'$, which is affine, by (c). Since $G/N$ is of finite type over $k$ (cf. VIA, 3.3.2), this implies, by EGA I, 5.1.10, that $G/N$ is affine.
% typo: est est affine -> est affine.

Finally, for arbitrary $k$, let $\overline k$ be an algebraic closure of $k$. Then, by 9.2 (v), $(G\otimes_k\overline k)/(N\otimes_k\overline k)$ is isomorphic to $(G/N)\otimes_k\overline k$, so, since the former is affine, so is the latter, and thus $G/N$ is also affine by (fpqc) descent (cf. EGA IV$_2$, 2.7.1). This proves 11.17 when $G$ is of finite type. To extend this to the general case, we shall need the following lemma.\nde{This lemma, taken from [DG70], \S\,III.3, 7.1, has been added.}

\begin{lemma}{11.17.1}\label{VIB.11.17.1}
Let $(C_i\to A_i)_{i\in I}$ be a filtered inductive system of ring morphisms, all faithfully flat. Then $A=\varinjlim A_i$ is faithfully flat over $C=\varinjlim C_i$.
\end{lemma}
\begin{proof}
By [BAC] \S\,I.3, Prop. 9, for a ring morphism $B\to B'$ to be faithfully flat, it is necessary and sufficient that it be injective and $B'/B$ be a flat $B$-module. Since each $C_i\to A_i$ is faithfully flat, one therefore has exact sequences
\[
0\to C_i\to A_i\to A_i/C_i\to0
\]
and $A_i/C_i$ is a flat $C_i$-module, so $(A_i/C_i)\otimes_{C_i}C$ is a flat $C$-module. Since inductive limits are exact and commute with tensor products, one obtains an exact sequence
\[
0\to C\to A\to A/C\to0
\]
as well as an isomorphism
\[
\varinjlim\bigl((A_i/C_i)\otimes_{C_i}C\bigr)=(A/C)\otimes_C C=A/C,
\]
from which one deduces that $A/C$ is a flat $C$-module. Thus $A$ is faithfully flat over $C$.
\end{proof}
Let us now return to the proof of 11.17 in the general case, i.e.\ when $G$ is not supposed of finite type. Put $\mathscr A(G)=A$ and $\mathscr A(N)=B=A/J$. By 11.13, $A$ is the inductive limit of an increasing filtered family $(A_i)_{i\in I}$ of Hopf subalgebras of finite type, so $G$ is the projective limit of the affine algebraic $k$-groups $G_i=\Spec(A_i)$. Denote by $\Delta$ (resp. $\tau$) the comultiplication (resp. antipode) of $A$, and $\Delta^2=(\Delta\otimes\mathrm{id}_A)\circ\Delta$.

For every $i$, $B_i=A_i/(J\cap A_i)$ is a quotient Hopf algebra of $A_i$, so $N_i=\Spec(B_i)$ is a closed subgroup of $G_i$. Moreover, since $N$ is invariant in $G$, the morphism $G\times N\to G$ defined by $(g,n)\mto gng^{-1}$ factors through $N$, and this is equivalent to saying that the pair $(A,J)$ satisfies the following property:
\[
\bigl(m_{13}\circ(\Delta\otimes\tau)\circ\Delta\bigr)(J)\subset A\otimes_k J
\]
where $m_{13}$ denotes the map $a_1\otimes a_2\otimes a_3\mto a_1a_3\otimes a_2$.
